\section{Closing discussion (Day 2, late afternoon)}

\subsection{Roundtable on Monte Carlo}
The final MC roundtable consolidated the actions identified on the Day-2
morning session: extend the COMETA polarized-$ZZ$ comparison to $WZ$ and to
semi-leptonic/boosted final states; converge on documented conventions for
polarized signal definitions (frames, DPA/NWA, interference) to be used by
ATLAS, CMS, and the generator authors alike; and quantify parton-shower and
hadronization uncertainties specifically for polarization-sensitive
substructure observables, where low-level taggers were shown to be most
vulnerable.

\subsection{Towards a COMETA deliverable}
The workshop closed with a planning session on the COMETA deliverable that
should crystallize the outcome of the meeting. The envisaged product combines:
(i) a common, public, polarization-labeled dataset of boosted hadronic
$W/Z$ jets and QCD background (following the open-data proposal of Day 1),
with agreed truth definitions grounded in the MC roundtable conventions;
(ii) a benchmark comparison of the tagging strategies presented---traditional
substructure, Particle-Transformer classification and regression, quantum
Lund-plane networks, decorrelated (3D~DDT) taggers, and amplitude-assisted
optimal observables---with common metrics including modeling robustness; and
(iii) a road map for the first hadronic/semi-leptonic polarization
measurements on the Run-2+3 datasets and their HL-LHC projections. A written
COMETA report collecting these elements, with contributions from the workshop
speakers, was agreed as the concrete next step.
